% Supplementary material. Requires booktabs, multirow, amsmath.
% Tables live in docs/tables/ and are generated by
%   scripts/supplement_tables.py   (tab_setup, tab_rayfan, tab_cluster, tab_ringrot)
%   scripts/paper_artifacts.py     (tab_main_multi and the main-paper tables)
% Do not edit a table by hand: the generators read the pipeline's own outputs,
% which is what keeps the prose and the numbers from drifting apart.

\appendix
\section{Supplementary material}
\label{sec:supp}

This appendix records what a reader would need in order to reproduce the
numbers and, more importantly, what we found while checking them. Three of the
four subsections exist because an assumption we had made turned out to be wrong
when measured, and we would rather state those plainly than leave them for a
reader to rediscover.

\subsection{Benchmarks: what is shared and what is not}
\label{sec:supp-setup}

% Generated by scripts/supplement_tables.py -- do not edit by hand.
\begin{table}[t]\centering\small
\caption{What the three benchmarks share and where they differ. The acoustic simulation is identical; the cameras are not. The field of view fixes how many rays of the matched fan fit inside the image (Tab.~\ref{tab:rayfan}), and Structured3D ships neither a furnished mesh nor a fixed camera height, so its query is rendered on the same floorplan proxy as the candidates and at a height that varies within a scene while the candidate grid uses one height per scene.}
\label{tab:setup}
\begin{tabular}{lccc}\toprule
 & Replica & Matterport3D & Structured3D \\\midrule
Sampling rate & 48000\,Hz & 48000\,Hz & 48000\,Hz \\
Microphones & 6 & 6 & 6 \\
Ring radius & 0.05\,m & 0.05\,m & 0.05\,m \\
DESDF cell / orientations & 0.1\,m / 36 & 0.1\,m / 36 & 0.1\,m / 36 \\
Depth rays / max range & [40, 160] / 20.0\,m & [40, 160] / 20.0\,m & [40, 160] / 20.0\,m \\
Image & $640{\times}480$ & $640{\times}480$ & $640{\times}360$ \\
Horizontal FOV & $106.2602^\circ$ & $106.2602^\circ$ & $80.0^\circ$ \\
$F/W$ & 0.3750 & 0.3750 & 0.5959 \\
Rays in the matched fan & 11 & 11 & 7 \\
Views per pose & 4 & 4 & 1 \\
Motion collections & forward, general & forward, general & none (fixed cameras) \\
Camera height & $1.25$\,m & $1.25$\,m & per frame, $0.36$\,m spread within a scene \\
Furnished query & yes & yes & no (no scan mesh) \\
\bottomrule\end{tabular}\end{table}


Tab.~\ref{tab:setup} is read from the dataset metadata rather than written by
hand. The acoustic simulation is identical on all three benchmarks: $48$\,kHz,
six omnidirectional receivers on a $5$\,cm ring with a co-located source,
reflection depth $50$, diffraction to order $10$, $20{,}000$ indirect rays. The
pose grid and the depth supervision are identical too.

The cameras are not, and three of the differences change what the pipeline has
to do.

\textbf{Field of view.} Replica and Matterport3D render at $106.26^\circ$;
Structured3D renders at $80^\circ$. This is not cosmetic, for the reason in
Sec.~\ref{sec:supp-fan}.

\textbf{Views per pose.} Replica and Matterport3D have navigable scan meshes, so
the simulator walks an agent and renders four-frame chunks at $0.27$\,m spacing,
in two motion regimes matching Gibson's forward and general sets.
Structured3D has no mesh: it ships perspective renders at the houses' own fixed
camera positions, one per pose, $0.8$--$1.1$\,m apart. Multi-view matching needs
translation between views to recover structure, so Structured3D admits only the
monocular backbone, and we report it that way rather than synthesising views we
would then be measuring instead of the baseline.

\textbf{Query acoustics.} Replica and Matterport3D ship a furnished scan, so
their query is what a microphone in the furnished room would record and the
floorplan candidates it is matched against are missing the furniture. That gap
is the problem this paper addresses. Structured3D has no furnished mesh at all,
so its query can only be rendered on the same floorplan proxy as the candidates.
Its rows are the matched-geometry condition, which is a different and much
easier problem, and they must not be read as a stronger result.

One smaller asymmetry works against Structured3D and we record it for
completeness: its camera height varies per frame, by a median of $0.36$\,m within
a single scene, while the candidate grid is rendered at one receiver height per
scene. Replica and Matterport3D fix the camera at $1.25$\,m and have no such
mismatch.

\subsection{The matched fan must fit inside the image}
\label{sec:supp-fan}

% Generated by scripts/supplement_tables.py -- do not edit by hand.
\begin{table}[t]\centering\small
\caption{The matched fan has to fit inside the image. Rays are $10^\circ$ apart, so $V$ rays span $(V{-}1)\cdot 10^\circ$, and the depth vector only exists out to the angle of its outermost column. Beyond it the interpolation returns NaN and localisation collapses without raising. Measured on three Structured3D scenes, $30$ queries.}
\label{tab:rayfan}
\begin{tabular}{ccccccc}\toprule
$V$ & $F/W$ & fan half-angle & image half-angle & $1$\,m & $2$\,m & median \\\midrule
11 & 0.5959 & $50^\circ$\,$\dagger$ & $39.3^\circ$ & 3.3 & 3.3 & 7.92\,m \\
9 & 0.5959 & $40^\circ$\,$\dagger$ & $39.3^\circ$ & 3.3 & 3.3 & 7.92\,m \\
7 & 0.5959 & $30^\circ$ & $39.3^\circ$ & 23.3 & 26.7 & 4.62\,m \\
5 & 0.5959 & $20^\circ$ & $39.3^\circ$ & 13.3 & 16.7 & 4.86\,m \\
11 & 0.375 & $50^\circ$ & $52.4^\circ$ & 20.0 & 20.0 & 5.65\,m \\
\bottomrule\end{tabular}
\\[2pt]\footnotesize $\dagger$ part of the fan falls outside the image.
\end{table}


The observation model compares $V$ rays, spaced $10^\circ$ apart, against $V$
consecutive orientation bins of the DESDF. The rays are read out of the
predicted depth vector by interpolating at image columns
$w = \tan\theta\, W\, F/W + (W{-}1)/2$. The depth vector only exists between its
first and last column, so a ray whose angle exceeds the image's own half-angle
interpolates outside the data and returns NaN. Nothing raises; recall simply
collapses.

At $106.26^\circ$ the outermost stored column sits at $52.5^\circ$ and the
eleven-ray fan, which spans $\pm 50^\circ$, fits. At $80^\circ$ the outermost
column sits at $39.3^\circ$: eleven rays put four outside the image and nine put
two outside, leaving seven as the largest odd fan that fits. Tab.~\ref{tab:rayfan}
shows what that costs in practice. Eleven rays at the Structured3D field of view
reach $3.3\%$ recall at $1$\,m; seven rays reach $23.3\%$. Keeping eleven rays
and pretending the camera is wider, $F/W = 0.375$, avoids the NaN but reads each
ray from the wrong column and reaches $20.0\%$.

We report this because the failure is silent and because the correct value is
not a tuned one: it follows from the field of view and the $10^\circ$ ray
spacing, and the sweep only confirms it.

\subsection{What the confidence intervals establish}
\label{sec:supp-stats}

% Generated by scripts/supplement_tables.py -- do not edit by hand.
\begin{table}[t]\centering\small
\caption{Two bootstraps of the same Replica result. Resampling queries answers whether the gain would survive more poses in these three rooms. Resampling rooms answers whether it would survive a new building, and it is the question the paper asks: queries inside a room share a floorplan, a candidate grid and an acoustic field, so with three rooms the effective sample size is three. We report the query interval in the main tables and this one beside it, and we do not claim building-level generalisation from Replica alone.}
\label{tab:cluster}
\begin{tabular}{lccccc}\toprule
Backbone & queries & rooms & $\Delta_{1\mathrm{m}}$ & query bootstrap & room bootstrap \\\midrule
F3Loc mono & 300 & 3 & +6.3 & $[+2.0, +10.7]$ & $[-3.0, +13.0]$ \\
UnLoc & 300 & 3 & +5.0 & $[+0.7, +9.3]$ & $[-3.0, +14.0]$ \\
DisCo-FLoc RRP & 300 & 3 & +10.0 & $[+6.0, +14.3]$ & $[+9.0, +12.0]$ \\
\bottomrule\end{tabular}\end{table}


The intervals in the main paper are paired bootstraps over queries. They answer
the question \emph{would this gain survive more poses in these rooms}, and they
are the right intervals for comparing two decision rules on one benchmark.

They are not intervals on building-level generalisation, and Tab.~\ref{tab:cluster}
makes the difference explicit. Queries inside a room share a floorplan, a
candidate grid and an acoustic field; they are not independent draws. Resampling
whole rooms instead of queries widens every Replica interval to span zero,
because Replica's test split has three rooms and the effective sample size is
three, not three hundred.

We therefore do not claim building-level generalisation from Replica. The
benchmarks that can support such a claim are the ones with many test rooms:
Structured3D contributes $28$ rooms and Matterport3D $12$, both held out by room
rather than by pose collection, and their intervals in Tab.~\ref{tab:main_multi}
are computed the same way.

\subsection{The motion collections are not a holdout}
\label{sec:supp-split}

Each benchmark ships two pose collections, forward and general motion, following
Gibson's convention. It is tempting to fit the four scalars on one and report on
the other, and an earlier version of this work did exactly that.

It does not hold anything out that matters here. On the Replica test rooms,
$73\%$ of general-motion poses have a forward-motion pose within $0.1$\,m, one
grid cell, and $100\%$ have one within the $1$\,m threshold being reported. The
acoustic score is a function of position alone, so a scalar fitted on the
forward collection has already seen the acoustic evidence of nearly every
reported position. The split holds out the visual side, which is not the side
being tuned.

Every number in this paper therefore pools the two collections and holds out
whole rooms instead: leave-one-room-out on Replica's three rooms, and alternate
rooms on the benchmarks with more. This doubles the test set and costs nothing
to render.

\subsection{The ring carries no heading}
\label{sec:supp-ring}

% Generated by scripts/supplement_tables.py -- do not edit by hand.
\begin{table}[t]\centering\small
\caption{The six-microphone ring carries no heading. The query's channels are rotated by $c$ positions and matched against an unrotated grid of $600$ candidates; the rank of the true cell barely moves. This is why the query ring rotating with the camera while the candidate ring is fixed to the world costs nothing, and it is also why every reported orientation comes from the visual backbone. A binaural pair is the change that would make the acoustic score a function of heading, at one engine call per (cell, heading).}
\label{tab:ringrot}
\begin{tabular}{ccc}\toprule
Channel rotation & mean rank of the true cell & worst of five \\\midrule
0 & 0.0 & 0 \\
1 & 0.0 & 0 \\
2 & 0.0 & 0 \\
3 & 0.6 & 2 \\
4 & 0.2 & 1 \\
5 & 0.0 & 0 \\
\bottomrule\end{tabular}\end{table}


The candidate grid stores one acoustic entry per cell, rendered with the ring
fixed to the world frame, while a query's ring rotates with the camera so that
channel $3$ is always camera-forward. These two conventions disagree, and
Tab.~\ref{tab:ringrot} shows that the disagreement costs nothing: rotating a
query's channels by any amount leaves the true cell first among $600$
candidates.

That is a convenience and a limitation at once. A $5$\,cm ring at $48$\,kHz gives
inter-microphone delays of about seven samples with a co-located source, so the
six channels are close to one channel repeated. It is the reason the acoustic
score is flat across the $36$ yaw bins and every reported orientation comes from
the visual backbone.

Making the acoustic term contribute to heading requires receivers that are
directional, such as a binaural pair with a head-related transfer function, and
a candidate grid rendered once per (cell, heading) rather than once per cell.
That is a factor of six more engine calls per cell at $36$ headings and, for the
three Replica test scenes, $86$\,GB of impulse responses. We report the
measurement rather than adopt the cost, since $88\%$ of the poses our method
places within $1$\,m are already within $30^\circ$, leaving little for a
heading-aware acoustic term to recover.
